\documentclass[11pt,a4paper]{article}
\usepackage[utf8]{inputenc}
\usepackage[T1]{fontenc}
\usepackage[spanish]{babel}

% --- Tipografía premium: Palatino ---
\usepackage{mathpazo}
\usepackage[scaled=0.88]{beramono} % Monospaced más elegante
\usepackage{microtype} % Ajustes microtypográficos (kerning, protrusion)

% --- Geometría y espaciado ---
\usepackage{geometry}
\geometry{a4paper, top=2.8cm, bottom=2.8cm, left=3.0cm, right=3.0cm}
\usepackage{setspace}
\setstretch{1.20}
\setlength{\parskip}{0.8em}

% --- Paquetes de soporte ---
\usepackage{amsmath}
\usepackage{amssymb}
\usepackage{graphicx}
\usepackage{booktabs}
\usepackage{xcolor}
\usepackage{fancyhdr}
\usepackage{hyperref}
\usepackage{longtable}
\usepackage{enumitem}
\usepackage{titlesec}
\usepackage{tcolorbox}
\tcbuselibrary{listings, skins, breakable}
\usepackage{listings}
\usepackage{tikz}

% --- Paleta de colores ---
\definecolor{ucvblue}{HTML}{003366}
\definecolor{ucvgold}{HTML}{C5960C}
\definecolor{secgray}{HTML}{444444}
\definecolor{linkblue}{HTML}{1A5276}
\definecolor{codegreen}{HTML}{0F766E}  % Teal 700 para comentarios
\definecolor{codegray}{HTML}{64748B}   % Slate 500
\definecolor{codepurple}{HTML}{B45309} % Amber 700 para strings
\definecolor{codebg}{HTML}{F8FAFC}     % Slate 50 para fondo
\definecolor{codeframe}{HTML}{E2E8F0}  % Slate 200 para bordes
\definecolor{codeheader}{HTML}{F1F5F9} % Slate 100 para cabeceras
\definecolor{lightbg}{HTML}{F8F8F5}

% --- Hipervínculos ---
\hypersetup{
    colorlinks=true,
    linkcolor=linkblue,
    filecolor=ucvblue,
    urlcolor=linkblue,
    pdftitle={StreamUCV - Informe Técnico de la Solución},
    pdfauthor={Escuela de Computación, UCV},
}

% --- Formato de secciones ---
\titleformat{\section}
    {\Large\bfseries\color{ucvblue}}
    {\thesection.}{0.6em}{}
    [\vspace{0.2em}{\color{codeframe}\titlerule[0.8pt]}]

\titleformat{\subsection}
    {\large\bfseries\color{secgray}}
    {\thesubsection}{0.5em}{}

\titleformat{\subsubsection}
    {\normalsize\bfseries\color{secgray}}
    {\thesubsubsection}{0.5em}{}

\titlespacing*{\section}{0pt}{2.4em}{1.2em}
\titlespacing*{\subsection}{0pt}{1.8em}{0.8em}
\titlespacing*{\subsubsection}{0pt}{1.4em}{0.6em}

% --- Listas con más aire y orden ---
\setlist[itemize]{topsep=6pt, itemsep=4pt, parsep=2pt, leftmargin=1.8em}
\setlist[enumerate]{topsep=6pt, itemsep=4pt, parsep=2pt, leftmargin=1.8em}

% --- Bloques de código ---
\newcommand{\macbuttons}{%
    \tikz[baseline=-0.6ex]{%
        \fill[color=red!70!white] (0,0) circle (2.2pt);%
        \fill[color=yellow!70!orange] (7pt,0) circle (2.2pt);%
        \fill[color=green!70!black] (14pt,0) circle (2.2pt);%
    }%
}

\lstdefinestyle{mystyle}{
    basicstyle=\ttfamily\footnotesize,
    commentstyle=\color{codegreen}\itshape,
    keywordstyle=\color{ucvblue}\bfseries,
    stringstyle=\color{codepurple},
    breakatwhitespace=false,
    breaklines=true,
    keepspaces=true,
    numbers=none,
    showspaces=false,
    showstringspaces=false,
    showtabs=false,
    tabsize=2,
}
\lstset{style=mystyle}

% Bloque de código con barra lateral de acento y esquinas redondeadas
\newtcblisting{codeblock}[1]{%
    enhanced,
    colback=codebg,
    colframe=codeframe,
    arc=5pt,
    boxrule=0.5pt,
    borderline west={3pt}{0pt}{ucvblue},
    left=10pt, right=10pt, top=8pt, bottom=8pt,
    listing only,
    listing options={style=mystyle, #1},
}

% Bloque de código con etiqueta de lenguaje y botones estilo macOS
\newtcblisting{codeblocklang}[2]{%
    enhanced,
    colback=codebg,
    colframe=codeframe,
    arc=5pt,
    boxrule=0.5pt,
    borderline west={3pt}{0pt}{ucvblue},
    left=10pt, right=10pt, top=8pt, bottom=8pt,
    coltitle=secgray,
    fonttitle=\small\ttfamily\bfseries,
    colbacktitle=codeheader,
    titlerule=0.5pt,
    title={\macbuttons\hfill\textnormal{\textit{#2}}},
    listing only,
    listing options={style=mystyle, #1},
}

% --- Encabezados y pies de página ---
\setlength{\headheight}{14pt}
\addtolength{\topmargin}{-2pt}
\pagestyle{fancy}
\fancyhf{}
\renewcommand{\headrulewidth}{0.4pt}
\fancyhead[L]{\small\color{secgray}\textit{StreamUCV --- Proyecto \#1}}
\fancyhead[R]{\small\color{secgray}\textit{Administración de Bases de Datos}}
\fancyfoot[R]{\small\color{codegray}Página \thepage}

\begin{document}

% ==============================================================================
% PORTADA
% ==============================================================================
\begin{titlepage}
    \centering
    \vspace*{0.5cm}
    
    % Logos de la Universidad y Facultad a cada lado del encabezado
    \begin{minipage}{0.18\textwidth}
        \flushleft
        \includegraphics[height=2.2cm]{../logo_central.png}
    \end{minipage}%
    \begin{minipage}{0.64\textwidth}
        \centering
        {\large \textbf{UNIVERSIDAD CENTRAL DE VENEZUELA}}\\[0.15cm]
        {\large \textbf{FACULTAD DE CIENCIAS}}\\[0.15cm]
        {\large \textbf{ESCUELA DE COMPUTACIÓN}}\\[0.15cm]
        {\large \textbf{ADMINISTRACIÓN DE BASES DE DATOS}}
    \end{minipage}%
    \begin{minipage}{0.18\textwidth}
        \flushright
        \includegraphics[height=2.2cm]{../logo_ciencias.png}
    \end{minipage}
    
    \vspace{5cm}
    
    {\Large \textbf{PROYECTO \#1 --- STREAMUCV}}\\[0.4cm]
    {\Large \textbf{INFORME TÉCNICO DE LA SOLUCIÓN}}\\[2.2cm]
    
    \vspace{5cm}

    \begin{minipage}{0.4\textwidth}
        \begin{flushleft} \large
            \emph{Grupo Docente:}\\
            Profa. Concettina Di Vasta \\
            Br. Joiner Rojas\\
            Br. Maximiliano Barboza
        \end{flushleft}
    \end{minipage}
    ~
    \begin{minipage}{0.4\textwidth}
        \begin{flushright} \large
            \emph{Integrantes:}\\
            Ronald Herrera\\
            Brhandon Palomo\\
            Sebastián Russian\\
        \end{flushright}
    \end{minipage}\\[3cm]
    
    {\large Caracas, Junio de 2026}
\end{titlepage}

\newpage
\tableofcontents
\newpage

% ==============================================================================
% SECCIÓN 1: DESCRIPCIÓN GENERAL DE LA SOLUCIÓN
% ==============================================================================
\section{Descripción General de la Solución del Proyecto}
\textbf{StreamUCV} es una aplicación en \textbf{Python} con interfaz \textbf{Streamlit} que consulta en tiempo real el \textbf{Diccionario de Datos} de una instancia \textbf{SQL Server 2022} y presenta de forma visual los metadatos del esquema \texttt{streaming}.

\subsection{Arquitectura de la Solución}
La solución se organiza en capas lógicas con responsabilidades bien delimitadas, asegurando modularidad y mantenimiento sencillo:
\begin{description}[leftmargin=4.2cm, labelwidth=3.9cm, style=sameline, font=\ttfamily\bfseries\color{ucvblue}]
    \item[main.py] Punto de entrada principal que inicializa Streamlit y lanza la aplicación.
    \item[app/config.py] Centraliza los parámetros de conexión (\texttt{SERVER}, \texttt{DATABASE}, etc.) y las constantes físicas de cálculo (páginas de 8 KB y velocidad de 17 MB/s).
    \item[app/connection.py] Administra el canal de comunicación con SQL Server mediante \texttt{pyodbc}, gestionando los intentos de conexión física.
    \item[app/queries/] Módulos individuales (\texttt{r01.py} a \texttt{r10.py}) que aíslan la lógica de cada requerimiento del diccionario de datos, retornando conjuntos de datos listos para consumir.
    \item[app/ui/streamlit\_app.py] Define la interfaz gráfica, organizando las pestañas, selectores dinámicos y la barra de navegación lateral.
\end{description}

\subsection{Beneficios del Diseño}
La arquitectura de la solución aporta ventajas clave frente a una ejecución SQL de consola tradicional:
\begin{description}[leftmargin=2.8cm, labelwidth=2.5cm, style=sameline, font=\bfseries\color{ucvblue}]
    \item[Portabilidad] Los parámetros de conexión pueden modificarse y probarse en tiempo de ejecución desde la interfaz web, permitiendo cambiar de servidor sin alterar el código.
    \item[Eficiencia] Las consultas se realizan sobre el Diccionario de Datos del motor (\texttt{sys.*}), evitando escaneos costosos sobre datos de usuario reales.
    \item[Interactividad] Transforma la administración pasiva de la base de datos en una experiencia visual enriquecida mediante tablas dinámicas, filtros reactivos y reportes limpios.
\end{description}

\newpage

% ==============================================================================
% SECCIÓN 2: EXPLICACIÓN DEL USO DEL DICCIONARIO DE DATOS
% ==============================================================================
\section{Explicación del Uso del Diccionario de Datos}
El Diccionario de Datos (D/D) es el conjunto de vistas internas del esquema \texttt{sys} que almacenan los metadatos de SQL Server. La aplicación consulta exclusivamente estas vistas para obtener información sobre la estructura lógica y física de la base de datos.

\subsection{R1: Listado de Tablas e Índices}
Lista las tablas e índices del esquema \texttt{streaming}.
\begin{itemize}
    \item \textbf{Vistas utilizadas}: \texttt{sys.tables}, \texttt{sys.schemas}, \texttt{sys.indexes}.
    \item \textbf{Lógica}: Se unen las vistas filtrando por \texttt{s.name = 'streaming'}. Para índices se excluyen las filas con \texttt{i.name IS NULL} (Heap).
\end{itemize}

\subsection{R2: Conteo de Tablas e Índices por Tabla}
Indica la cantidad total de tablas y cuántos índices tiene cada una.
\begin{itemize}
    \item \textbf{Vistas utilizadas}: \texttt{sys.tables}, \texttt{sys.schemas}, \texttt{sys.indexes}.
    \item \textbf{Lógica}: Conteo escalar de tablas y \texttt{LEFT JOIN} con \texttt{sys.indexes} agrupando por tabla.
\end{itemize}

\subsection{R3: Restricciones del Esquema}
Enumera las restricciones de integridad declaradas.
\begin{itemize}
    \item \textbf{Vistas utilizadas}: \texttt{sys.objects}, \texttt{sys.tables}, \texttt{sys.schemas}.
    \item \textbf{Clasificación}: Se filtran por \texttt{c.type}:
        \begin{center}
            \small
            \begin{tabular}{ll}
                \toprule
                \textbf{Código} & \textbf{Tipo de Restricción} \\
                \midrule
                \texttt{PK} & Clave Primaria \\
                \texttt{F}  & Clave Foránea \\
                \texttt{C}  & Check Constraint \\
                \texttt{UQ} & Unicidad \\
                \texttt{D}  & Valor por Defecto \\
                \bottomrule
            \end{tabular}
        \end{center}
\end{itemize}

\subsection{R4: Columnas de cada Índice}
Detalla qué columnas componen cada índice y su orden.
\begin{itemize}
    \item \textbf{Vistas utilizadas}: \texttt{sys.indexes}, \texttt{sys.index\_columns}, \texttt{sys.columns}.
    \item \textbf{Lógica}: Se usa \texttt{STRING\_AGG} ordenando por \texttt{ic.key\_ordinal}, indicando \texttt{ASC}/\texttt{DESC} según \texttt{ic.is\_descending\_key}.
\end{itemize}

\subsection{R5: Triggers del Esquema}
Lista los disparadores asociados a las tablas.
\begin{itemize}
    \item \textbf{Vistas utilizadas}: \texttt{sys.triggers}, \texttt{sys.trigger\_events}, \texttt{sys.tables}.
    \item \textbf{Lógica}: Se clasifica el tipo (\texttt{AFTER} o \texttt{INSTEAD OF}) y el estado (habilitado/deshabilitado). Los eventos se agrupan con \texttt{STRING\_AGG}.
\end{itemize}

\subsection{R6: Tamaño Físico por Tabla}
Muestra el espacio en disco ocupado por cada tabla.
\begin{itemize}
    \item \textbf{Vistas utilizadas}: \texttt{sys.tables}, \texttt{sys.dm\_db\_partition\_stats}.
    \item \textbf{Fórmulas de Almacenamiento} (SQL Server almacena datos en páginas de $8\text{ KB}$):
        \begin{equation}
            \text{Espacio de Datos (KB)} = \text{Páginas de Datos} \times 8
        \end{equation}
        \begin{equation}
            \text{Espacio de Índices (KB)} = \text{Páginas de Índices} \times 8
        \end{equation}
        Donde las páginas de datos corresponden a registros con \texttt{index\_id} $\in \{0,1\}$ (tabla base/Heap o índice clúster), y las páginas de índices a registros con \texttt{index\_id} $> 1$ (índices no agrupados).
\end{itemize}

\subsection{R7: Tamaño Estimado del Registro}
Estima cuántos bytes ocupa un registro completo de cada tabla en promedio.
\begin{itemize}
    \item \textbf{Vistas utilizadas}: \texttt{sys.columns}, \texttt{sys.tables}.
    \item \textbf{Fórmula de Estimación}:
        \begin{equation}
            \text{Tamaño del Registro} = \sum \left( \text{Longitud de cada Columna} \right)
        \end{equation}
        Se obtiene sumando los tamaños máximos (\texttt{max\_length}) de todas las columnas de la tabla.
        
        \textit{Ejemplo práctico:} Una columna \texttt{VARCHAR(50)} aporta $50\text{ bytes}$, mientras que una de tipo \texttt{NVARCHAR(50)} aporta $100\text{ bytes}$ (debido a la codificación de doble byte de caracteres Unicode).
\end{itemize}

\subsection{R8: Tamaño de cada Columna}
Desglosa el tipo de dato y tamaño físico de cada columna.
\begin{itemize}
    \item \textbf{Vistas utilizadas}: \texttt{sys.columns}, \texttt{sys.types}.
    \item \textbf{Lógica}: Se reconstruye la definición SQL de la columna a partir de su tipo mediante sentencias condicionales (\texttt{varchar(n)}, \texttt{nvarchar(n/2)}, \texttt{decimal(p,s)}, etc.).
\end{itemize}

\subsection{R9: Factor de Bloqueo}
El \textbf{Factor de Bloqueo} indica cuántos registros caben como máximo en una única página física de 8 KB ($8192\text{ bytes}$).
\begin{itemize}
    \item \textbf{Factor de Bloqueo de la Tabla}:
        \begin{equation}
            \text{Factor de Bloqueo}_{\text{Tabla}} = \left\lfloor \frac{8192}{\text{Tamaño del Registro}} \right\rfloor
        \end{equation}
    \item \textbf{Factor de Bloqueo del Índice} (basado solo en las columnas indexadas):
        \begin{equation}
            \text{Factor de Bloqueo}_{\text{Índice}} = \left\lfloor \frac{8192}{\text{Tamaño de la Clave}} \right\rfloor
        \end{equation}
        Donde el \textit{Tamaño de la Clave} es la suma de las longitudes de las columnas que forman parte de dicho índice.
\end{itemize}

\subsection{R10: Costo de Consultas de Igualdad}
Estima el costo de ejecutar una consulta de tipo \texttt{WHERE columna = valor} en términos de accesos a páginas y tiempo.

\subsubsection{1. Detección de Índices}
Un índice es utilizable para búsquedas directas si la columna por la que se consulta es la primera clave registrada en la estructura (\texttt{key\_ordinal = 1}).

\subsubsection{2. Caso A: Escaneo Completo (Full Table Scan)}
Sin un índice útil, SQL Server debe leer todas las páginas físicas de la tabla:
\begin{equation}
    \text{Costo}_{\text{Scan}} = \text{Total de Páginas}
\end{equation}
\begin{equation}
    \text{Tiempo}_{\text{Scan}} = \frac{\text{Total de Páginas} \times 8\text{ KB}}{17\text{ MB/s}}
\end{equation}

\subsubsection{3. Caso B: Búsqueda por Índice (Index Seek)}
Con un índice aplicable, se recorre su estructura en árbol B-Tree. La altura aproximada del árbol ($h$) se estima con un factor de ramificación (\textit{fan-out}) promedio de $100$:
\begin{equation}
    \text{Altura del Árbol } (h) = \max\left(1, \left\lceil \log_{100}(\text{Total de Páginas}) \right\rceil\right)
\end{equation}

El costo se calcula dependiendo del tipo de índice detectado:
\begin{itemize}
    \item \textbf{Índice Clustered:} $\text{Costo}_{\text{Seek}} = h + 1$ (las hojas contienen directamente los datos).
    \item \textbf{Índice Non-Clustered:} $\text{Costo}_{\text{Seek}} = h + 2$ (requiere el recorrido y luego leer la página de datos o \textit{lookup}).
\end{itemize}

Para tablas con pocas páginas, se limita el costo del Seek para que no supere al de un Scan completo:
\begin{equation}
    \text{Costo Seek Final} = \min(\text{Costo}_{\text{Seek}},\; \text{Costo}_{\text{Scan}})
\end{equation}

El tiempo estimado de búsqueda física se deduce homólogamente:
\begin{equation}
    \text{Tiempo}_{\text{Seek}} = \frac{\text{Costo Seek Final} \times 8\text{ KB}}{17\text{ MB/s}}
\end{equation}

\newpage

% ==============================================================================
% SECCIÓN 3: INSTRUCCIONES DE EJECUCIÓN
% ==============================================================================
\section{Instrucciones de Ejecución}

\subsection{Paso 1: Prerrequisitos}
\begin{enumerate}
    \item \textbf{Python 3.11+}. Verificar con:
    \begin{codeblock}{language=bash}
python --version
    \end{codeblock}
    \item \textbf{Microsoft SQL Server} activo y accesible.
    \item \textbf{Driver ODBC}: \texttt{ODBC Driver 17 for SQL Server} (recomendado) u \texttt{ODBC Driver 18}.
\end{enumerate}

\subsection{Paso 2: Instalar Dependencias}
Desde la raíz del proyecto (\texttt{Fase-1}):
\begin{codeblocklang}{language=bash}{Terminal}
pip install -r requirements.txt
\end{codeblocklang}
Esto instala \textbf{Streamlit} (interfaz web), \textbf{pyodbc} (conexión ODBC) y \textbf{pandas} (manejo de datos).

\subsection{Paso 3: Crear y Poblar la Base de Datos}
El script \texttt{setup\_db.py} ejecuta en orden los archivos SQL de la carpeta \texttt{scripts/}:
\begin{enumerate}
    \item \texttt{create\_repositorio\_sqlserver.sql}: Crea la BD \texttt{StreamUCV} y el esquema \texttt{streaming}.
    \item \texttt{create\_tables\_sqlserver.sql}: Crea tablas, claves, restricciones e índices.
    \item \texttt{insert\_tables\_sqlserver.sql}: Carga los datos de prueba.
\end{enumerate}

\begin{codeblocklang}{language=bash}{Terminal}
python setup_db.py
\end{codeblocklang}
\textit{Nota: Si falla la conexión, verifique que \texttt{SERVER} en \texttt{app/config.py} coincida con su instancia (ej. \texttt{localhost\textbackslash SQLEXPRESS}).}

\subsection{Paso 4: Configurar la Conexión}
Edite \texttt{app/config.py} según su entorno:
\begin{codeblocklang}{language=python}{app/config.py}
SERVER = "localhost\\SQLEXPRESS"
DATABASE = "StreamUCV"
USERNAME = ""       # Vacio para Windows Auth
PASSWORD = ""
DRIVER = "ODBC Driver 17 for SQL Server"
\end{codeblocklang}
Si \texttt{USERNAME} y \texttt{PASSWORD} quedan vacíos, se usa \texttt{Trusted\_Connection=yes} (autenticación de Windows).

\subsection{Paso 5: Ejecutar la Aplicación}
\begin{codeblocklang}{language=bash}{Terminal}
streamlit run main.py
\end{codeblocklang}
Se abrirá automáticamente en \url{http://localhost:8501}.

\subsection{Paso 6: Navegación}
\begin{enumerate}
    \item \textbf{Barra Lateral}: Permite modificar los parámetros de conexión en caliente y probar la conexión sin reiniciar.
    \item \textbf{Pestañas principales}:
    \begin{itemize}
        \item \textbf{Estructura Lógica}: Tablas e índices (R1), conteos (R2), restricciones (R3), detalle de índices (R4).
        \item \textbf{Disparadores}: Triggers del esquema (R5).
        \item \textbf{Almacenamiento}: Tamaño por tabla (R6), tamaño de registro (R7), tamaño por columna (R8).
        \item \textbf{Factor de Bloqueo}: Registros por página de tablas e índices (R9).
        \item \textbf{Costo de Consulta}: Selección de tabla/columna para comparar Full Scan vs. Index Seek (R10).
    \end{itemize}
\end{enumerate}

\end{document}
